\section{Introduction}

Picture two drones surveying a road. The first reports a blockage. The second,
looking from a different angle, reports the same thing. A task manager that
counts distinct \texttt{source\_id} values sees two sources and closes the
task. That step is safe only if the two reports can fail separately. When both
drones look at the same parked truck, the same shadow, or the same low-contrast
vehicle, they can miss together, and agreement carries less evidence than the
count suggests.

Classifier-ensemble research has measured this problem for decades. Kuncheva
and Whitaker~\cite{kuncheva2003} study the double-fault measure, the rate at
which two classifiers are wrong on the same sample, and compare it with the
rate expected if the classifiers erred independently. We port that statistic
into collaborative visual perception, where the ``samples'' are physical
objects and the ``classifiers'' are agents that may differ in model, position,
and view. We do not propose a new metric. The contribution is a measured
benchmark, a calibration procedure, and an accounting of what naive agreement
costs in task decisions.

We ask three questions, each answered only for the corpus we built.
\begin{enumerate}
  \item How far above the independence prediction is the joint miss rate of
        two agents, and does it differ by detector pairing?
  \item Can that inflation be fitted on labelled scenes and carried to scenes
        and illuminations the fit never saw?
  \item What does a closing rule that assumes independence cost in missed
        hazards, and what does a more cautious rule cost in verifications?
\end{enumerate}

\paragraph{Scope.} Everything here is simulation: CARLA 0.9.16, one town, two
detectors, vehicle objects, one frame per scene. We make no claim about field
deployments or about detectors we did not run.

\paragraph{Contributions.}
(i)~A balanced 900-scene corpus with per-viewpoint, pixel-derived visibility
truth. (ii)~Inflation ratios with bootstrap intervals for three detector
pairings, with a held-out-illumination check. (iii)~A decision-level comparison
of two closing rules, reported together with its verification cost and with the
limit that the cautious rule did not use a dependence estimate.
